\begin{figure}[htbp]
\centering
\begin{tikzpicture}[x=1.05cm,y=1cm,>=Latex,font=\small]
 \draw[very thick,ink] (0,0)--(9,0);
 \draw[very thick,dashed,accent] (9,0)--(10,0);
 \foreach \x in {0,...,10}{
   \draw (\x,-.08)--(\x,.08);
   \node[below=5pt,font=\scriptsize] at (\x,0) {$\x$};
 }
 \foreach \x in {1,...,9}{
  \draw[->,ink] (\x-.85,.35)--(\x-.15,.35);
  \node[above=3pt] at (\x-.5,.35) {$\x$};
 }
 \draw[->,accent] (9.15,.35)--(9.85,.35);
 \node[above=3pt,text=accent] at (9.5,.35) {$1$};
 \draw[decorate,decoration={brace,amplitude=5pt},ink] (0,1)--(9,1)
  node[midway,above=7pt] {Nueve intervalos orientados: una vuelta};
 \node[align=center] at (5,-1.05)
  {Índice de tránsito: $g(t)=1+(t\bmod9)$\\
   Memoria del calendario levantado: $m(t)=\lfloor t/9\rfloor$};
 \node[align=center] at (5,-1.85)
  {$g(t+9)=g(t),\qquad m(t+9)=m(t)+1$};
\end{tikzpicture}
\caption{Levantamiento del ciclo de nueve tránsitos. Los números superiores
etiquetan intervalos orientados; los inferiores localizan sus fronteras.
El intervalo discontinuo muestra el primer tránsito de la vuelta siguiente.
El índice de fase retorna y el cociente del calendario aumenta. Este esquema
representa el calendario; la acción sobre residuos, cocientes, orientación,
frontera y memoria se define mediante las operaciones del TPK.}
\label{fig:ii-segmento-generador}
\end{figure}
